% ============================= VSR results =============================
\section{Orientation Is the Persistent Weak Family}
\label{sec:vsr}

\subsection{Main results}
Table~\ref{tab:vsr} shows overall and orientation accuracy for all eight
conditions. Orientation is the weakest family for every condition by a wide
margin: $62.8\%$ (2B zero-shot) and $63.5\%$ (7B zero-shot) against
$75$--$93\%$ for the easy families (e.g., containment $89.5\%$, vertical
$85.2\%$). Scaling 2B$\rightarrow$7B raises overall accuracy by $6.9$ points but
orientation by $0.7$ points. LM-side LoRA is the only intervention that moves
orientation above $64\%$ ($65.7\%$), and hard-negative training adds a further
$0.7$ points ($66.4\%$); both are far below every other family. Within the
family, \emph{facing away from} is the hardest relation ($48.7\%$ zero-shot,
$59.0\%$ best) and is the relation where hard negatives help most
($59.0\%\rightarrow66.7\%$); \emph{perpendicular to} ($n{=}12$) never exceeds
$58.3\%$ and degrades under LM-side training.

\begin{table*}[t]
\centering
\small
\setlength{\tabcolsep}{3.4pt}
\begin{threeparttable}
\caption{\textbf{Main VSR results (test, $n{=}2{,}195$; orientation $n{=}137$).}
Every value is recomputed from the committed prediction CSVs and matches the
published metrics JSONs exactly. Per-relation cells: facing $n{=}64$, facing away
$n{=}39$, parallel $n{=}22$, perpendicular $n{=}12$.}
\label{tab:vsr}
\begin{tabular}{lcc|cccc}
\toprule
Condition & Overall & Orient. & Facing & Away & Parall. & Perp. \\
\midrule
2B zero-shot       & 74.0 & 62.8 & 70.3 & 53.8 & 59.1 & 58.3 \\
2B General LoRA    & 76.6 & 62.0 & 68.8 & 53.8 & 59.1 & 58.3 \\
2B Targeted LoRA   & 76.5 & 64.2 & 73.4 & 53.8 & 54.5 & 66.7 \\
2B structured prompt & 68.3\tnote{a} & --- & --- & --- & --- & --- \\
\midrule
7B zero-shot       & 80.9 & 63.5 & 73.4 & 48.7 & 63.6 & 58.3 \\
7B General LoRA    & \textbf{84.7} & \textbf{65.7} & \textbf{75.0} & 59.0 & 63.6 & 41.7 \\
7B HardNeg LoRA    & 84.3 & \textbf{66.4} & 70.3 & \textbf{66.7} & \textbf{68.2} & 41.7 \\
7B Projector LoRA  & 82.9 & 64.2 & 78.1 & 56.4 & 50.0 & 41.7 \\
7B Vision+Proj LoRA & 83.1 & 64.2 & 76.6 & 51.3 & 63.6 & 41.7 \\
\bottomrule
\end{tabular}
\begin{tablenotes}\footnotesize
\item[a] Single-run structured-prompt variant (2B); orientation cells not
isolated for this condition.
\end{tablenotes}
\end{threeparttable}
\end{table*}

\begin{figure}[t]
\centering
\includegraphics[width=\columnwidth]{fig/orientation_per_relation.pdf}
\caption{\textbf{Orientation per-relation accuracy} across the four headline
conditions ($n$ annotated above each bar). Hard-negative training is the only
condition that lifts \emph{facing away from} above $60\%$, at a cost to
\emph{facing}; no condition fixes \emph{perpendicular to}.}
\label{fig:perrel}
\end{figure}

\subsection{The failure is reasoning-shaped, not perception-shaped}
We annotated the 48 test statements that fail persistently --- wrong under
both 7B zero-shot and 7B LoRA, or wrong across $\geq$3 of the 4 base
conditions, or wrong under both 2B and 7B zero-shot --- with an
eight-class failure taxonomy (App.~F). The distribution is
decisive for locating the failure:
\textbf{37.5\% are clear-image reasoning failures} (the pose is obvious to a
human, the model still errs), 16.7\% are genuine viewpoint ambiguity, 12.5\%
angular-geometry failures, 10.4\% have questionable ground-truth labels, and only
\textbf{4.2\% are vision-limited} (small/occluded objects). ``Facing'' /
``facing away'' statements make up 68.8\% of persistent failures, and in 9 cases
the 7B model \emph{regresses} on clear cases that the 2B model answers correctly.
Two implications: (i) the label noise ceiling is real ($\sim$10\% questionable
labels) but cannot explain the family-level gap; we confirm this below with a
clean-subset evaluation; (ii) the residual is not at the photoreceptor level, so
we test the representation and behavior hypotheses explicitly in the next two
sections.

\subsection{Clean-subset control}
On the 48 annotated failures, we removed (a) the 13 with questionable labels or
genuine ambiguity (``clear'' subset, $n{=}124$) and (b) a further 10 with
ambiguous intrinsic orientation or small/occluded objects (``strict'' subset,
$n{=}114$). Probe and generative accuracies barely move on either subset
(e.g., the best ungrounded probe: $68.9\% \rightarrow 71.1\% \rightarrow 71.6\%$;
generative LoRA T1: $68.9\% \rightarrow 72.2\% \rightarrow 78.4\%$); the
orientation deficit is not a label-noise artifact (full table in
App.~C).
